% Teil II: der durchgerechnete Weg.

% Gedankenschritt: Kürzen + Symmetrie; "mindestens" = Summe;
% Säulen-Formel (Kugeln + Trennstriche).

\subsection*{Gegeben und gesucht}
(a) $\binom{12}{3}$, $\binom{57}{56}$, $\binom{10}{7}$; (b) Anzahl der
Möglichkeiten, aus $10$ Personen mindestens $7$ auszuwählen; (c) Anzahl der
Lösungen von $x+y+z=19$ mit $x,y,z\in\mathbb N_0$.

\subsection*{a) Werte}
\begin{align}
  \binom{12}{3}&=\frac{12\cdot11\cdot10}{3\cdot2\cdot1}=220,\\
  \binom{57}{56}&=\binom{57}{1}=57, \quad\text{(Symmetrie: wähle die
        ``eine, die nicht in der Teilmenge ist``)}\\
  \binom{10}{7}&=\binom{10}{3}=\frac{10\cdot9\cdot8}{6}=120 .
\end{align}

\subsection*{b) Mindestens 7 aus 10}
``Mindestens 7`` schließt 7, 8, 9, 10 ein. Summe über $k=7,\dots,10$:
\begin{equation}
  \sum_{k=7}^{10}\binom{10}{k}=120+45+10+1=176 .
\end{equation}
\textbf{Gedanke:} ``Mindestens`` ist eine \emph{Summe} über alle
jenseitsliegenden $k$; der Fehler ``nur $\binom{10}{7}$`` würde die drei
Grob-Mengen auslassen.

\subsection*{c) Zahl der Lösungen $x+y+z=19$, $\mathbb N_0$}
Säulen-Formel mit $N=19$, $d=3$ (drei Unbekannte):
\begin{equation}
  \binom{19+3-1}{3-1}=\binom{21}{2}=\frac{21\cdot20}{2}=210 .
\end{equation}
\textbf{Warum $\binom{21}{2}$?} Stelle 19 ``Kugeln`` (die Einsen) und 2
``Trennstriche`` (zwischen $x,y,z$) in einer Reihe von 21 Symbolen an;
jede Wahl der 2 Strich-Positionen aus 21 gibt eine Lösung
$x+y+z=19$ mit $x,y,z\ge0$.

\begin{ergebnis}
\begin{equation*}
  \text{a)}\ 220,\ 57,\ 120,\qquad
  \text{b)}\ 176,\qquad
  \text{c)}\ 210 .
\end{equation*}
\end{ergebnis}

\begin{kontrolle}
\textbf{a) Nachrechnen:} $\binom{10}{8}=45$, $\binom{10}{9}=10$,
$\binom{10}{10}=1$, Summe $120+45+10+1=176$. \checkmark

\textbf{c) Plausibilität:} $\binom{21}{2}=210$; alternativ zählt man
$y=0,\dots,19$ durch: für festes $y$ ist $x+z=19-y$ mit $20-y$ Lösungen;
$\sum_{y=0}^{19}(20-y)=20+19+\dots+1=210$. Konsistent.

\textbf{Grenzfall:} W\"aren $x,y,z$\emph{strikt} positiv (alle $\ge1$),
substituiert man $x'=x-1$ usw. und erhält $x'+y'+z'=16$ mit $x'$
$\ge0$: $\binom{16+3-1}{2}=\binom{18}{2}=153$. Der Vergleich zeigt, wie
die Verschiebung die Zahl senkt (weniger ``Boden``-Lösungen).
\end{kontrolle}

\begin{merksatz}
$\binom{n}{k}$ z\"ahlt ungeordnete Auswahlmengen (Kern: Reihenfolge
``wegdividieren`` durch $k!$); ``mindestens`` ist eine Summe, und
nichtnegative Lösungen einer Summengleichung z\"ahlt die
S\"aulen-Formel $\binom{N+d-1}{d-1}$.
\end{merksatz}
